\documentclass[10pt]{article}

\usepackage[utf8]{inputenc}

\usepackage[T1]{fontenc}

\usepackage{amsmath}

\usepackage{amsfonts}

\usepackage{amssymb}

\usepackage[version=4]{mhchem}

\usepackage{stmaryrd}

\usepackage{bbold}

\usepackage{graphicx}

\usepackage[export]{adjustbox}

\graphicspath{ {./images/} }



\begin{document}

относительно группы локальных замен координат в $\left(\mathbb{R}^{m}, 0\right)$ может быть взята деформация



\[

F(\lambda)=f+\sum_{i=1}^{\mu} \lambda_{i} \varphi_{i},

\]



где $\varphi_{i}, i=1,2, \ldots, \mu,-$ ростки функций, образуюшие базис в факторкольце кольца ростков функций на $\left(\mathbb{R}^{m}, 0\right)$ по идеалу, порожденному частными производными $\partial f / \partial x_{i}$ функции $f$ (якобиеву идеалу функции $f$ ). Если в семействе общего положения не встречаются элементы, имеюшие при классификации с точностью до эквивалентностей из группы $G$ непрерывные модули, то это семейство является версальным для каждого своего элемента. Напр., это имеет место для зависящего не более чем от 5 параметров семейства гладких функций на произвольном многообразии.



О. В. Ляшко.



В. Д. квадратичного гамильтониана $H$, заданного на конечномерном векторном симплектич. пространстве $V$ - содержащее $H$ семейство $F(x)$ квадратичных гамильтонианов на $V$ (где $x \in \mathbb{R}^{k}, F(x)$ гладко зависит от $x$ и $F(0)=H$, обладающее следующим свойством: лля любого квадратичного гамильтониана $G$ на $V$, близкого к $H$, существуют точка $x=x_{G} \in \mathbb{R}^{k}$ и линейная канонич. замена координат $\boldsymbol{\Phi}=\boldsymbol{\Phi}_{G}$ в $V$, переводящая $G$ в $F\left(x_{G}\right)$, причем $x_{G}$ и $\boldsymbol{\Phi}_{G}$ гладко зависят от $G, x_{H}=0$ и $\Phi_{H}$ тождественна. Имеются явные выражения для В. д. на вещественном пространстве $V$ (см. [1], [2]).



Если $\lambda_{1}, \lambda_{2}, \ldots-$ попарно различные собственные числа линейного гамильтонова оператора, соответствующего данному квадратичному гамильтониану $H$, а $m_{1}\left(\lambda_{j}\right) \geqslant m_{2}\left(\lambda_{j}\right) \geqslant$ $\geqslant m_{3}\left(\lambda_{j}\right) \geqslant \ldots-$ порядки жордановых клеток, отвечающих $\lambda_{j}$, то минимально возможное количество $k$ параметров В. д. $H$ равно



\[

\frac{1}{2} \sum_{j}\left(m_{1}\left(\lambda_{j}\right)+3 m_{2}\left(\lambda_{j}\right)+5 m_{3}\left(\lambda_{j}\right)+\ldots\right)+\frac{n}{2},

\]



где $n-$ количество жордановых клеток нечетных порядков с собственным числом 0 (см. [2]).



Лит.: [1] Гал и н Д. М., «Труды семинара имени И. Г. Петровского», 1975, т. 1, с. 63-74; [2] K ос ak H., «J. Diff. Equat.», 1984, v. 51, № 3, p. 359-407.



ВЕРСАЛЬНОЕ СЕМЕЙСТВО - сМ. БИфУркацИЯ.



BEPTИКАЛЬНЫЙ BEKTOP - см. Голономии группа. ВЕРШИННАЯ ДИАТРАММА ФЕЙНМАНА - сМ. Фейнмана диаграмма.



BEРШИННАЯ ЧАСТЬ, в е рши н ная ф функция,- одна из основных функций в квантовой теории поля, характеризующая взаимодействие между квантовыми полями; содержит все радиационные поправки. В перенормированной теории возмущений В.ч. определяется как сумма вкладов, отвечающих сильно связным Фейнмана диаграммам с числом и типом внешних линий, определяемых соответствующей вершиной в правилах Фейнмана.



\begin{center}

\includegraphics[max width=\textwidth]{2023_07_08_3b1333c734b5ab5977d7g-1}

\end{center}



Вктады однопетлевых $\left(\alpha \Gamma_{\mu}^{(1)}\right)$ и двухпетлевых $\left(\alpha^{2} \Gamma_{\mu}^{(2)}\right)$ диаграмм в вершинной части $\Gamma_{\mu}\left(p, p^{\prime} ; q\right)$ в квантовой электродинамике: $p, p^{\prime}$ и $\boldsymbol{q}$ - соответственно 4-импульсы начального и конечного электрона и фотона, $\gamma_{\mu}$ - матрицы Дирака $(\mu=0,1,2,3)$. Так, напр., В. ч. $\Gamma_{\mu}\left(p, p^{\prime} ; q\right)$ в квантовой электродинамике определяется суммой (перенормированных) вкладов, к-рые, по правилам Фейнмана, изображаются диаграммами (рис.) и представляются в виде степенного ряда по безразмерному параметру (постоянной тонкой структуры) $\alpha \approx 1 / 137$, характеризующему интенсивность электромагнитных процессов.



При более общем определении, без обращения к теории возмушений, вершинные части (как и полные Грина функции) выражаются через вариационные производные от производящего функционала.



Иногда (неправильно) В. ч. называют просто в е рш и Н о й.



Д. В. Ширков.



BEC а в томорфной формы - см. Автоморфная форма. BEC модуля рной фо рмы - см. Модулярная форма.



BEC представления $\rho$ алгебры Ли в векторн о м п р остран с т в е - отображение $\alpha$ алгебры Ли $L$ в ее поле определения $k$, для которого существует такой ненулевой вектор $x$ пространства $V$, что



\[

(\rho(h)-\alpha(h) 1)^{n_{x, h}}(x)=0

\]



для всех $h \in L$ и некоторого целого $n_{x, h}>0$ (вообще говоря, зависящего от $x$ и $h$ ), где 1 обозначает тождественное преобразование $V$. В этом случае говорят также, что $\alpha-$ в е с $L$-м о дуля $V$, определяемого представлением $\rho$. Множество всех векторов $x \in V$, удовлетворяющих указанному условию, вместе с нулем образует подпространство $V_{\alpha}$, называемое весовым подпространством веса $\alpha$ (или, соответствующим весу $\alpha$ ). Если $V=V_{\alpha}$, то $V$ называется весовым пространством, или весовым моду ле м над $L$, веса $\alpha$.



Если $V$ и $W-$ весовые модули над $L$ весов $\alpha$ и $\beta$ соответственно, то их тензорное произведение $V \otimes W$ является весовым модулем веса $\alpha+\beta$. Если $L-$ нильпотентная алгебра Ли, то весовое подпространство $V_{\alpha}$ веса $\alpha$ в $V$ является подмодулем $L$-модуля $V$. Если, кроме того,



\[

\operatorname{dim}_{k} V<\infty,

\]



a $\rho(L)$ - расщепляемая алгебра Ли линейных преобразований модуля $V$, то $V$ разлагается в прямую сумму конечного числа весовых $L$-подмодулей разных весов:



\[

V=V_{\sigma} \oplus V_{\delta} \oplus \ldots \oplus V_{\tau}

\]



(весовое разложение $V$ относительно $L$ ). Если $L-$ нильпотентная подалгебра конечномерной алгебры Ли $M$, рассматриваемой как $L$-модуль относительно присоединенного представления $\mathrm{ad}_{M}$ алгебры $M$, и ad $_{M} L$ является расщепляемой алгеброй Ли линейных преобразований $M$, то соответствующее весовое разложение $M$ относительно $L$ :



\[

M=M_{\alpha} \oplus M_{\beta} \oplus \ldots \oplus M_{\gamma}

\]



называется разложением Фиттинга $M$ относительно $L$, веса $\alpha, \beta, \ldots, \gamma$ называются к о р н я м и, а пространства $M_{\alpha}$, $M_{\beta}, \ldots, M_{\gamma}-$ корневыми подпространствами $M$ относительно $L$. Если, кроме того, задано представление $\rho$ алгебры $M$ в конечномерном векторном пространстве $V$, для К-рого $\rho(L)$ - расщепляемая алгебра Ли линейных преобразований $V$, и



\[

V=V_{\sigma} \oplus V_{\delta} \oplus \ldots \oplus V_{\tau}

\]



\begin{itemize}

  \item соответствующее весовое разложение $V$ относительно $L$, To

\end{itemize}



\[

\rho\left(M_{\alpha}\right)\left(V_{\sigma}\right) \subseteq V_{\alpha+\sigma},

\]



когда $\alpha+\sigma$ есть вес $V$ относительно $L$, и



\[

\rho\left(M_{\alpha}\right)\left(V_{\sigma}\right)=0

\]



в противном случае. В частности, если $\alpha+\beta-$ корень, то



\[

\left[M_{\alpha}, M_{\beta}\right] \subseteq M_{\alpha+\beta},

\]



в остальных случаях



\[

\left[M_{\alpha}, M_{\beta}\right]=0 .

\]





\end{document}